%%%%%%%%%%%%%%%%%%%%%%%%%%%%%%%%%%%%%%%%%%%%%%%%%%%%%%%%%%%%%%%%%%%%%%%%%%%%%%%
% Full instructions available at:
% https://github.com/elauksap/focus-beamertheme

\documentclass{beamer}
\usetheme{focus}

% package for font
\usepackage{fontspec}
\defaultfontfeatures{Mapping=tex-text}  %%如果没有它，会有一些 tex 特殊字符无法正常使用，比如连字符。
\usepackage{xunicode,xltxtra}
\usepackage[BoldFont,SlantFont,CJKnumber,CJKchecksingle]{xeCJK}  % \CJKnumber{12345}: 一万二千三百四十五
\usepackage{CJKfntef}  %%实现对汉字加点、下划线等。
\usepackage{pifont}  % \ding{}
% package for math
\usepackage{amsfonts}
% package for graphics
\usepackage[americaninductors,europeanresistors]{circuitikz}
\usepackage{tikz}
\usetikzlibrary{plotmarks}  % placements=positioning
\usepackage{graphicx}  % \includegraphics[]{}
\usepackage{subfigure}  %%图形或表格并排排列
% package for table
\usepackage{colortbl,dcolumn}  %% 彩色表格
\usepackage{multirow}
\usepackage{multicol}
\usepackage{booktabs}
% package for code
\usepackage{fancyvrb}
\usepackage{listings}
% setting for font
%%\setCJKmainfont{Adobe Kaiti Std}
\setCJKmainfont{SimSun}
%%\setCJKmainfont{Microsoft YaHei}
%%\setCJKmainfont{FangSong_GB2312} 
%%\setmainfont{Times New Roman}
%\setsansfont[Mapping=tex-text]{Adobe Song Std}
%如果装了Adobe Acrobat，可在font.conf中配置Adobe字体的路径以使用其中文字体。
%也可直接使用系统中的中文字体如SimSun、SimHei、微软雅黑等。
%原来beamer用的字体是sans family；注意Mapping的大小写，不能写错。
%设置字体时也可以直接用字体名，以下三种方式等同：
%\setromanfont[BoldFont={黑体}]{宋体}
%\setromanfont[BoldFont={SimHei}]{SimSun}
%\setromanfont[BoldFont={"[simhei.ttf]"}]{"[simsun.ttc]"}

% font setting by xeCJK
\setCJKfamilyfont{NSimSun}{NSimSun}
\newcommand{\song}{\CJKfamily{NSimSun}}
%%%\setCJKfamilyfont{AdobeSongStd}{Adobe Song Std}
%%%\newcommand{\AdobeSong}{\CJKfamily{AdobeSongStd}}
\setCJKfamilyfont{FangSong}{FangSong_GB2312}
\newcommand{\fang}{\CJKfamily{FangSong}}
%%%\setCJKfamilyfont{AdobeFangsongStd}{Adobe Fangsong Std}
%%%\newcommand{\AdobeFang}{\CJKfamily{AdobeFangsongStd}}
\setCJKfamilyfont{SimHei}{SimHei}
\newcommand{\hei}{\CJKfamily{SimHei}}
%%%\setCJKfamilyfont{AdobeHeitiStd}{Adobe Heiti Std}
%%%\newcommand{\AdobeHei}{\CJKfamily{AdobeHeitiStd}}
\setCJKfamilyfont{KaiTi}{KaiTi}
\newcommand{\kai}{\CJKfamily{KaiTi}}
%%%\setCJKfamilyfont{AdobeKaitiStd}{Adobe Kaiti Std}
\newcommand{\AdobeKai}{\CJKfamily{AdobeKaitiStd}}
\setCJKfamilyfont{LiSu}{LiSu}
\newcommand{\li}{\CJKfamily{LiSu}}
\setCJKfamilyfont{YouYuan}{YouYuan}
\newcommand{\you}{\CJKfamily{YouYuan}}
\setCJKfamilyfont{FZJingLei}{FZJingLeiS-R-GB}
\newcommand{\jinglei}{\CJKfamily{FZJingLei}}
\setCJKfamilyfont{MSYH}{Microsoft YaHei}
\newcommand{\msyh}{\CJKfamily{MSYH}}

\setbeamerfont{frametitle}{series=\Large\hei} % 修改Beamer标题字体

\graphicspath{{figures/}}  %%图片路径
\renewcommand\figurename{图}

% setting for pdf
\hypersetup{% pdfpagemode=FullScreen,%
  pdfauthor={Xiaodong Wang},%
  pdftitle={Title},%
  CJKbookmarks=true,%
  bookmarksnumbered=true,%
  bookmarksopen=false,%
  plainpages=false,%
  colorlinks=true,%
  citecolor=green,%
  filecolor=magenta,%
  linkcolor=blue,%red(default)
  urlcolor=cyan}

% setting for fontspec
\XeTeXlinebreaklocale "zh"  %%表示用中文的断行
\XeTeXlinebreakskip = 0pt plus 1pt minus 0.1pt  %%多一点调整的空间
%%%%%

% 自定义颜色
\def\Red{\color{red}}
\def\Green{\color{green}}
\def\Blue{\color{blue}}
\def\Mage{\color{magenta}}
\def\Cyan{\color{cyan}}
\def\Brown{\color{brown}}
\def\White{\color{white}}
\def\Black{\color{black}}

\lstnewenvironment{xmlCode}[1][]{% for Java
  \lstset{
    basicstyle=\tiny\ttfamily,%
    columns=flexible,%
    framexleftmargin=.7mm, %
    % frame=shadowbox,%
    % rulesepcolor=\color{cyan},%
     frame=single,%
    backgroundcolor=\color{white},%
    xleftmargin=4\fboxsep,%
    xrightmargin=4\fboxsep,%
    numbers=left,numberstyle=\tiny,%
    numberblanklines=false,numbersep=7pt,%
    language=xml, %
    }\lstset{#1}}{}

\lstnewenvironment{javaCode}[1][]{% for Java
  \lstset{
    basicstyle=\tiny\ttfamily,%
    columns=flexible,%
    framexleftmargin=.7mm, %
    frame=shadowbox,%
    rulesepcolor=\color{cyan},%
    % frame=single,%
    backgroundcolor=\color{white},%
    xleftmargin=4\fboxsep,%
    xrightmargin=4\fboxsep,%
    numbers=left,numberstyle=\tiny,%
    numberblanklines=false,numbersep=7pt,%
    language=Java, %
    }\lstset{#1}}{}

\lstnewenvironment{shCode}[1][]{% for Java
  \lstset{
    basicstyle=\scriptsize\ttfamily,%
    columns=flexible,%
    framexleftmargin=.7mm, %
    frame=shadowbox,%
    rulesepcolor=\color{brown},%
    % frame=single,%
    backgroundcolor=\color{white},%
    xleftmargin=4\fboxsep,%
    xrightmargin=4\fboxsep,%
    numbers=left,numberstyle=\tiny,%
    numberblanklines=false,numbersep=7pt,%
    language=sh, %
    }\lstset{#1}}{}

\newcommand\ask[1]{\vskip 4bp \tikz \node[rectangle,rounded corners,minimum size=6mm,
  fill=white,]{\Cyan \includegraphics[height=1.5cm]{question} \Large \msyh #1};}

\newcommand\wxd[1]{\vskip 4bp \tikz \node[rectangle,minimum size=6mm,
  fill=blue!60!white,]{\White \ding{118} \msyh #1};}

\newcommand\cxf[1]{\vskip 4bp \tikz \node[rectangle,rounded corners,minimum size=6mm,
  fill=purple!60!white,]{\White \ding{42} \msyh #1};}

\newcommand\xyy[1]{\vskip 2bp \tikz \node[rectangle,minimum size=3mm,
  fill=black!80!white,]{\White \msyh\scriptsize #1};}

\newcommand\tta[1]{\vskip 4bp \tikz \node[rectangle,minimum size=6mm,
  fill=blue!60!white,]{\White \ding{118} \msyh #1};}

\newcommand\ttb[1]{\vskip 4bp \tikz \node[rectangle,rounded corners,minimum size=6mm,
  fill=purple!60!white,]{\White \ding{42} \msyh #1};}

\newcommand\ttc[1]{\vskip 2bp \tikz \node[rectangle,minimum size=3mm,
  fill=black!80!white,]{\White \msyh\scriptsize #1};}

\newcommand\homework[1]{\vskip 2bp \tikz \node[rectangle,minimum size=3mm,
  fill=red!80!white,]{\White \ding{45} \msyh\scriptsize 课后小作业 } ; {\kai\small #1}} 

\newcommand\notice[1]{\vskip 4bp \tikz \node[rectangle,rounded corners,minimum size=6mm,
  fill=red!80!white,]{\White \scriptsize \ding{42} \msyh #1};}

\newcommand\samp[1]{\vskip 2bp \tikz \node[rectangle,minimum size=3mm,]{\Mage\msyh \small CODE \ding{231} \Black #1};\vskip -8bp}

\newcommand\codeset[1]{\vskip 2bp \tikz \node[rectangle,minimum
  size=3mm]{\Mage\msyh \small 课程配套代码 \ding{231} \Black
    #1};\vskip -4bp}

\newcommand\pptlink[2]{\vskip 4bp \tikz \node[rectangle,rounded corners,minimum size=6mm,
  fill=blue!70!white,]{\href{run:#1}{\White \scriptsize \msyh 动画演示 #2}};}


\makeatletter
\newcommand{\Extend}[5]{\ext@arrow 0099{\arrowfill@#1#2#3}{#4}{#5}}
\makeatother
%%%%%%%%%%%%%%%%%%%%%%%%%%%%%%%%%%%%%%%%%%%%%%%%%%%%%%%%%%%%%%%%%%%%%%%%%%%%%%%
\title{\hei JAVA应用与开发\\  控制台应用程序设计}
\subtitle{\it 让我们愉快的Coding起来吧...}
\author{王晓东}
\titlegraphic{\vspace{-8em}\includegraphics[height=10cm]{static/ouc.pdf}\vspace{-6em}} 
\institute{中国海洋大学信息学院计算机系}
\date{\today}

\begin{document}
\begin{frame}
  \maketitle
\end{frame}

\begin{frame}
  \frametitle{学习目标}
  \begin{itemize}
  \item 了解计算机人机交互发展
  \item 掌握控制台程序设计开发中命令行参数、系统属性、标准输入输出的概
    念和相关Java操作
  \item 掌握Java文件操作的的常用方法
  \item 了解注解类型
  \item 学会Jar归档工具，包括通过命令行或IDE进行Java程序归档的方法
  \end{itemize}
\end{frame}

\begin{frame}
  \frametitle{大纲}
  \tableofcontents
\end{frame}

\section{从古老的计算机谈起}

\begin{frame}
  \frametitle{冯诺依曼机}

  \begin{columns}
    \column{0.5\textwidth}

    \ttc{1950年代 IBM 1401}
    \begin{figure}
      \centering
      \includegraphics[width=\textwidth]{fig-old-computer-01.jpg}
    \end{figure}

    \column{0.5\textwidth}
    
    \ttc{2010年代 树莓派开发板}

        \begin{figure}
      \centering
      \includegraphics[width=\textwidth]{fig-Raspberry-Pi-Board.png}
    \end{figure}
  \end{columns}

  \pause
  
  {\kai 我们的计算机是台遵守存储程序原理的冯诺依曼机器，基本组成包
    括{\Red\hei 运算器、控制器（合起来是CPU）、存储器、输入设备、输出设
      备}。你所面对的一切SOC也好，单板电脑也好，都是高度集成在一起的冯
    诺依曼机。}
\end{frame}

\begin{frame}
  \frametitle{人机交互}

  \tta{使用打孔卡片作为输入源，使用打印机作为输出设备}
  
  \begin{columns}
    \column{0.5\textwidth}

    一摞打孔卡片，就是一个“文件”。它可以是一段程序，也可以是一段程序需要使用的数据。
    
    \column{0.5\textwidth}
    \begin{figure}
      \centering
      \includegraphics[width=\textwidth]{fig-old-computer-input.jpg}
    \end{figure}
  \end{columns}

\end{frame}


\begin{frame}
  \frametitle{人机交互}

  \tta{BASIC语言解释器}

  纸带在70年代还很流行，当年比尔盖茨的BASIC语言解释器，就是存在纸带上的，现在已经成文物了。
  
  \begin{figure}
    \centering
    \includegraphics[width=0.8\textwidth]{fig-basic-lang-punch-tape.jpg}
  \end{figure}

\end{frame}

\begin{frame}
  \frametitle{人机交互}

  \tta{使用键盘作为输入设备，使用显示器作为输出设备}
  
  \begin{columns}
    \column{0.7\textwidth}

    \begin{figure}
      \centering
      \includegraphics[width=\textwidth]{fig-apple-1-modules.jpg}
    \end{figure}
    
    \column{0.3\textwidth}
    \begin{figure}
      \centering
      \includegraphics[width=\textwidth]{fig-apple-1.jpg}
    \end{figure}
  \end{columns}

\end{frame}

\begin{frame}
  \frametitle{人机交互}

  {\hei 再厉害的科幻片导演，在飞船的人机交互界面表达上也未能超越同时代计算机的发展。}
  

  \begin{figure}
    \centering
    \includegraphics[width=\textwidth]{fig-Al012.png}
  \end{figure}

\end{frame}


\section{命令行参数}

\begin{frame}[fragile] % [fragile]参数使得能够插入代码
\frametitle{命令行参数}

\begin{block}{命令行参数}
  在启动时Java控制台应用程序，可以一次性地向程序中传递（零至多个）字符串参
  数，这些参数被称为命令行参数。语法格式如下：

  \begin{shCode}
    java <应用程序类名> [<命令行参数>]*
  \end{shCode}
\end{block}

\pause

\tta{说明}

\begin{itemize}\kai
\item 命令行参数将被系统接收并静态初始化为一个一维的String数组对象，然后将之作为实参传给
  应用程序入口方法main()。
\item 命令行参数须使用空格符分隔，如果参数中包含空格符则必须使用双引号括起来。
\end{itemize}
\end{frame}

\begin{frame}[fragile] % [fragile]参数使得能够插入代码
  \frametitle{命令行参数}

  \codeset{sample.commandline.CommandLineArgsSample.java}

  \pause

  \ttc{Linux下运行程序}

  \begin{shCode}
    > java CommandLineArgsSample Lisa "Billy" "Mr Brown"
  \end{shCode}

  \pause

  \ttc{Windows下运行程序}

  \begin{shCode}
    C:\> java.exe CommandLineArgsSample Lisa "Billy" "Mr Brown" "a""b" 
  \end{shCode}

  \pause

  输出结果：

  \begin{stdoutCode}
    Lisa
    Billy
    Mr Brown
  \end{stdoutCode}
\end{frame}

\begin{frame}[fragile] % [fragile]参数使得能够插入代码
  \frametitle{可变参数方法}

  \begin{itemize}
  \item Java语言允许在定义方法时指定使用任意数量的参数，其格式是在参数
    类型后加“...”。
  \item 可变长度参数必须放在参数列表的最后，一个方法最多只能包含一个可
    变长度参数。
  \item 编译时，可变参数被当作{\Red\hei 一维数组处理}。
  \end{itemize}

  \begin{javaCode}
    public void myprint(String s, int i, Object... objs) { // 可变参数方法
      System.out.println(s.toUpperCase()); 
      System.out.println(100 * i);
      for(Object o: objs) { // 作为一维数组处理
        System.out.println(o); 
      }
    } 
  \end{javaCode}
\end{frame}

\section{系统属性}

\begin{frame}[fragile] % [fragile]参数使得能够插入代码
  \frametitle{系统属性}

  \begin{block}{Java系统属性}
    \begin{itemize}
    \item 记录当前操作系统和JVM等相关的环境信息。
    \item 以{\hei\Red 键值对}的形式存在，由{\hei\Red 属性名称、属性值}两部分组成。
    \item 均为字符串形式。
    \end{itemize}
  \end{block}

  \pause
  
  \tta{系统属性的用途}

  系统属性在URL网络编程、数据库编程和Java Mail邮件收发等编程中经常使用，
  一般被用来设置代理服务器、指定数据库的驱动程序类等。

  除了使用代码方法外，也可使用命令在运行程序时添加新的系统属性：
  
  \begin{shCode}
    >java -Dmmmm=vvvv SystemPropertiesSample
  \end{shCode}
  
\end{frame}


\begin{frame}[fragile] % [fragile]参数使得能够插入代码
  \frametitle{遍历、操作系统属性}

  \begin{block}{}
    可以使用System.getProperties()获得一个封装了当前运行环境下所有系统属
    性信息的Properties类（java.utils.Properties）的实例。
  \end{block}

  \codeset{sample.commandline.SystemPropertiesSample.java}

  \tta{可用方法}

  \begin{description}
  \item [\fbox{Enumeration propertyNames()}] %返回以Enumeration类型表示的所有可用系统属性的名称。
  \item [\fbox{String getProperty(String key)}] %获得特定系统属性的属性值。
  \item [\fbox{Object setProperty(string key, String value)}] %设置/添加单个系统属性信息。
  \item [\fbox{void load(InputStream inStream)}] 
  \item [\fbox{void store(OutputStream out, String header) }] %实现属性信息的导入/导出操作。
  \end{description}


\end{frame}

\section{标准输入/输出}

\begin{frame}[fragile] % [fragile]参数使得能够插入代码
  \frametitle{标准输入/输出}

  \tta{控制台程序的交互方式}

  \begin{itemize}
  \item 用户使用键盘作为{\hei 标准输入设备}向程序输入数据
  \item 程序利用计算机终端窗口作为{\hei 程序标准输出设备}显示输出数据
  \end{itemize}

  这种操作被称为{\hei 标准输入/输出（Standard Input/Output）}。

\end{frame}

\begin{frame}[fragile] % [fragile]参数使得能够插入代码
  \frametitle{标准输入/输出的分类}

  \begin{block}{}
    java.lang.System类的三个静态类成员提供了有关标准输入/输出的IO操作功能。
  \end{block}

  \begin{description}\kai
  \item [System.in] 从“标准输入”读入数据（java.io.InputStream类型）
  \item [System.out] 向“标准输出”写出数据（java.io.PrintStream类型）
  \item [System.err] 向“标准错误”写出数据（java.io.PrintStream类型）
  \end{description}

  \pause
  
  \tta{PrintStream类的主要方法}
  
  \begin{itemize}
  \item print()/println()方法 被进行了多次重载（boolean、char、int、 long、float、
    double以及char[], Object和String）。
  \end{itemize}
\end{frame}

\begin{frame}[fragile] % [fragile]参数使得能够插入代码
  \frametitle{读取控制台输入的传统方法}

  \begin{javaCode}
    import java.io.InputStreamReader; 
    import java.io.BufferedReader; 
    import java.io.IOException;

    public class TestStandardInput {
      public static void main (String args[]) {
        String s;
        InputStreamReader isr = new InputStreamReader(System.in); 
        BufferedReader br = new BufferedReader(isr);
        try {
          s = br.readLine(); 
          while (!s.equals("")) {
            System.out.println("Read: " + s);
            s = br.readLine(); 
          }
          br.close();
        } catch (IOException e) {
          e.printStackTrace(); 
        }
      } 
    }  
  \end{javaCode}
\end{frame}

\begin{frame}[fragile] % [fragile]参数使得能够插入代码
  \frametitle{对上述程序的几点解释}

  \begin{itemize}[<+-| structure@+>]
  \item System.in为InputStream类型对象，功能较弱，只能以字节为单位从预定义的标准输入（键盘）
    读取信息。
  \item 程序并没有直接操作System.in对象进行读取操作，而是将其封装为一个功能稍强
    的InputStreamReader对象，以字符为单位读取信息。\only<2>{实际的过程为：{\kai InputStreamReader对象并没
        有直接读取键盘输入，而是多次调用System.in对象的读字节功能，再将所得字节转换为字符。}}
  \item InputStreamReader仍不能令人满意，再次封装，得到BufferedReader对象。\only<3>{后者提供了缓冲读
      取的功能，即多次调用InputStreamReader读字符操作，然后将所读取的多个字符积累起来组成字符
      串，其间以换行符为分隔，最终实现以行为单位读取字符串功能。}
  \item 当在键盘上空回车时，BufferedReader的readLine()方法接收到的不是空值null，而是一个长
    度为零的字符串""，其中包含0个字符但仍然是一个Java对象。
  \end{itemize}
\end{frame}

\section{文件操作}

\begin{frame}[fragile] % [fragile]参数使得能够插入代码
  \frametitle{文件操作对象}

  \tta{创建File类对象}

  java.io包中定义与数据输入、输出功能有关的类，包括提供文件操作功能的File类。
  
  \ttc{public File(String pathname)}

  通过给定的路径/文件名字符串创建一个新File实例。
  
  \ttc{public File(String parent, String child)}

  通过分别给定的parent路径名和child文件名（也可以是子路径名）或字符串来创建一个新File实例。
\end{frame}

\begin{frame}[fragile] % [fragile]参数使得能够插入代码
  \frametitle{使用File类}

  \codeset{sample.commandline.FileOperationSample.java}

\end{frame}

\begin{frame}[fragile] % [fragile]参数使得能够插入代码
  \frametitle{File类的主要方法~\ding{182}} 

  \begin{columns}
    \column{0.5\textwidth}

    \tta{文件/目录名操作}

    \begin{itemize}
    \item String getName() 
    \item String getPath() 
    \item String getAbsolutePath() 
    \item String getParent()
    \end{itemize}

    \column{0.5\textwidth}

    \tta{设置和修改操作} 
    
    \begin{itemize}
    \item boolean delete() 
    \item void deleteOnExit() 
    \item boolean createNewFile() 
    \item setReadOnly() 
    \item boolean renameTo(File dest)
    \end{itemize}
  \end{columns}
\end{frame}

\begin{frame}[fragile] % [fragile]参数使得能够插入代码
  \frametitle{File类的主要方法~\ding{183}}

  \begin{columns}
    \column{0.5\textwidth}

    \tta{测试操作} 

    \begin{itemize}
    \item boolean exists() 
    \item boolean canWrite() 
    \item boolean canRead() 
    \item boolean isFile() 
    \item boolean isDirectory() 
    \item boolean isAbsolute()
    \end{itemize}

    \column{0.5\textwidth}

    \tta{目录操作}

    \begin{itemize}
    \item boolean mkdir() 
    \item String[] list() 
    \item File[] listFiles()
    \end{itemize}

    \tta{获取常规文件信息操作}

    \begin{itemize}
    \item long lastModified() 
    \item long length()
    \end{itemize}

  \end{columns}
\end{frame}


\begin{frame}[fragile] % [fragile]参数使得能够插入代码
  \frametitle{文件I/O有关读写类}

  \tta{常见的文本文件I/O操作的类}
  
  \ttc{java.io.FileReader类}
  
  提供read()方法以字符为单位从文件中读入数据。

  \ttc{java.io.FileWrite类}

  提供write()方法以字符为单位向文件写出数据。

  \ttc{java.io.BufferedReader类}

  提供readLine()方法以行为单位读入一行字符。

  \ttc{java.io.PrintWriter类}

  提供print()和println()方法以行为单位写出数据。

\end{frame}

\begin{frame}[fragile] % [fragile]参数使得能够插入代码
  \frametitle{读取文件内容}

  \samp{ReadFile.java}

  \begin{javaCode}
    import java.io.*;

    public class ReadFileSample {
      public static void main (String[] args) {
        String fname = "test.txt"; 
        File f = new File(fname);

        try {
          FileReader fr = new FileReader(f);  // 1
          BufferedReader br = new BufferedReader(fr); 
          String s = br.readLine();
          while (s != null) { // 2
            System.out.println("读入：" + s);
            s = br.readLine(); }
          br.close();
        } catch (FileNotFoundException e1) {
          System.err.println("File not found: " + fname); 
        } catch (IOException e2) {
          e2.printStackTrace(); 
        }
      } 
    }
  \end{javaCode}
\end{frame}

\begin{frame}[fragile] % [fragile]参数使得能够插入代码
  \frametitle{读取文件内容}

  \notice{上述代码几点说明}

  \begin{enumerate}
  \item FileReader的构造方法被重载过，接受以字符串形式给出的文件名，上述代码等价于：
    \begin{javaCode}
      FileReader fr = new FileReader("test.txt");
    \end{javaCode}
  \item 使用BufferedReader的readLine()方法读文件，遇到文件结尾则返回null，而不是""，与读取
    键盘输入遇到空回车时返回空字符串的情况不同。
  \end{enumerate}
  
\end{frame}

\begin{frame}[fragile] % [fragile]参数使得能够插入代码
  \frametitle{输出内容到文件}

  \samp{WriteFileSample}
  
  \begin{javaCode}
    import java.io.*;
    
    public class TestWriteFile {
      public static void main (String[] args) {
        File file = new File("tt.txt");
        try {
          InputStreamReader is = new InputStreamReader(System.in); 
          BufferedReader in=new BufferedReader(is);
          FileWriter fw = new FileWriter(file);
          PrintWriter out = new PrintWriter(fw);
          String s = in.readLine();
          while(!s.equals("")) { // 从键盘逐行读入数据输出到文件
            out.println(s);
            s = in.readLine(); 
          }
          in.close(); // 关闭 BufferedReader 输入流
          out.close(); // 关闭连接文件的 PrintWriter 输出流
        } catch (IOException e) {
          e.printStackTrace(); 
        }
      } 
    }
  \end{javaCode}
\end{frame}

\begin{frame}[fragile] % [fragile]参数使得能够插入代码
  \frametitle{对上述代码的几点说明}

  \begin{enumerate}
  \item 写文件时如果目标文件不存在，程序运行不会出错，而是自动创建该文
    件，但如果目标路径不存在，则会出错。
  \item 写文件操作结束后一定要关闭输出流，即关闭文件，否则被操作文件仍
    处于打开状态，不安全。
  \end{enumerate}
\end{frame}

\begin{frame}[fragile] % [fragile]参数使得能够插入代码
  \frametitle{文件过滤}

  \begin{block}{}
    文件过滤，即只检索和处理符合特定条件的文件。{\kai 最常见的为按照文
      件类型（后缀）进行划分，如查找.class或.xml文件。}

    文件过滤可以使用java.io.FileFilter接口，该接口只定义了一个抽象方
    法accept。
  \end{block}

  \ttc{boolean accept(File pathname)}

  测试参数指定的File对象对应的文件（目录）是否应该保留在文件列表中，即不被过滤。

  {\kai 在实际应用中，可以定义该接口的一个实现类，重写其中的accept()方法，在
  方法中添加文件过滤逻辑，然后创建一个该实现类的对象作为参数传递
  给File对象的文件列表方法list()，在list()方法执行过程中会自动调用前者
  的accept()方法来过滤文件。}
  
\end{frame}

\begin{frame}[fragile] % [fragile]参数使得能够插入代码
  \frametitle{使用FileFilter实现文件过滤}

  \codeset{sample.commandline.filefilter}
  
\end{frame}


%%%\section{过时API}
%%%\begin{frame}[fragile] % [fragile]参数使得能够插入代码
%%%  \frametitle{过时API}
%%%
%%%  过时API是指那些过去定义的，现已不提倡使用的API，包括类、属性和方法等。过时API均存在相应的
%%%  替代物,这些替代者可能采用了更标准化的命名惯例，或者功能更适用。在将来的JDK版本中，过
%%%  时API可能不再被支持，所以开发中应尽量避免使用。
%%%
%%%  \begin{javaCode}
%%%    import java.util.*;
%%%
%%%    public class TestDeprecation {
%%%      public static void main(String[] args) {
%%%        Date now = new Date();
%%%        int hour = now.getHours(); // 过时API
%%%        System.out.println(hour);
%%%      } 
%%%    }
%%%  \end{javaCode}
%%%
%%%
%%%\end{frame}
%%%
%%%\begin{frame}[fragile] % [fragile]参数使得能够插入代码
%%%  \frametitle{过时API}
%%%
%%%  编译程序时输出提示信息：
%%%  \begin{stdoutCode}
%%%    注意：TestDeprecation.java使用或覆盖了已过时的API。 
%%%    注意：要了解详细信息，请使用 -Xlint:deprecation 重新编译。
%%%  \end{stdoutCode}
%%%
%%%  使用下述命令重新编译程序：
%%%  \begin{shCode}
%%%    >javac -Xlint:deprecation TestDeprecation.java  
%%%  \end{shCode}
%%%
%%%  输出更详细说明信息：
%%%
%%%  \begin{stdoutCode}
%%%    TestDeprecation.java:5: 警告:[deprecation] java.util.Date 中的 
%%%    getHours() 已过时  
%%%    int hour = now.getHours(); ^
%%%    1警告
%%%  \end{stdoutCode}
%%%\end{frame}
%%%
%%%\begin{frame}[fragile] % [fragile]参数使得能够插入代码
%%%  \frametitle{对上述代码的改造}
%%%
%%%  在Java API文档中，java.util.Date类的getHour()部分已作如下说明：
%%%
%%%  {\kai “从JDK 1.1开始，由Calendar.get(Calendar.HOUR\_OF\_DAY)取代”}
%%%  \samp{TestDeprecation.java}
%%%  \begin{javaCode}
%%%    import java.utils.*;
%%%
%%%    public class TestDeprecation {
%%%      public static void main(String[] args) {
%%%        Calendar c = Calendar.getInstance();
%%%        int hour = c.get(Calendar.HOUR_OF_DAY);
%%%        System.out.println(hour);
%%%      }
%%%    }  
%%%  \end{javaCode}
%%%\end{frame}

\section{注解（Annotation）}

\begin{frame}[fragile] % [fragile]参数使得能够插入代码
  \frametitle{什么是注解}

  \begin{block}{注解}
    是从JDK5.0开始新添加的一种语言特性，区别于代码注释（Comment）。
  \end{block}

  \begin{itemize}[<+-|alert@+>]
  \item 注解不直接影响程序的语义，开发和部署工具可以对其读取并以某种形
    式处理这些注解，可能生成其他Java源文件、XML文档或要与包含注解的程序
    一起使用的其他构件。
  \item 本质上，注解就是可以添加到代码中的一种类似于修饰符的成分，可以
    用于声明包、类、构造方法、方法、属性、参数和变量等场合。
  \end{itemize}
\end{frame}

\begin{frame}[fragile] % [fragile]参数使得能够插入代码
  \frametitle{什么是注解}

  Java语言采用了一类新的数据类型来描述注解。（{\hei\Red 注解类型}）相当
  于类或接口，每一条注解相当于该注解类的一个实例。注解类型采用@interface标记来声明。

  \tta{JDK5.0及后续版本定义的几种有用的注解类型}
  
  \begin{itemize}
  \item public @interface Deprecated
  \item public @interface Override
  \item public @interface SuppressWarnings
  \end{itemize}
\end{frame}

\begin{frame}[fragile] % [fragile]参数使得能够插入代码
  \frametitle{Override注解}

  java.lang.Override类型注解用于指明被注解的方法重写了父类中的方法，如
  果不是合法的方法重写，则编译报错。

  \begin{javaCode}
    public class Person {
      ...
      @Override
      public String toString() { // 重写方法
        return "Name: " + name; 
      }
    }
  \end{javaCode}

  \ttc{toString的原始定义}
  
  \begin{javaCode}
    public String toString() {
      return getClass().getName() + "@" + Integer.toHexString(hashCode());
    }  
  \end{javaCode}
\end{frame}

\begin{frame}[fragile] % [fragile]参数使得能够插入代码
  \frametitle{Deprecated注解}

  Deprecated注解的作用是标记过时的API。如果通过方法重写或调用的方式来使
  用已被注解为过时的方法时，编译器将会根据注解信息发现不应该使用此方法，
  并作提醒。

  \begin{javaCode}
    public class A { 
      @deprecated
      public void ma() {
        System.out.println("In class A, just for test!");
      } 
    }
  \end{javaCode}

\end{frame}

\begin{frame}[fragile] % [fragile]参数使得能够插入代码
  \frametitle{SuppressWarnings注解}

  使用SuppressWarnings注解可以关闭编译器对指定的一种或多种问题的提示/警
  告功能。

  \tta{语法格式，比较自由，下述均可} 

  \begin{javaCode}
    @SuppressWarnings(value={"deprecation"})
    @SuppressWarnings(value={"deprecation","unchecked"}) 
    @SuppressWarnings("deprecation") 
    @SuppressWarnings({"deprecation", "unchecked"})
  \end{javaCode}

  \begin{javaCode}
    import java.util.*;
    import java.lang.SuppressWarnings; 

    @SuppressWarnings(value={"deprecation"}) 
    public class TestSuppressWarnings {
      public static void main(String[] args) {
        Date now = new Date();
        int hour = now.getHours(); 
        System.out.println(hour);
      } 
    }  
  \end{javaCode}

  代码编译时，则不会再输出先前的提示API过时信息。
\end{frame}

\section{归档工具}

\begin{frame}[fragile] % [fragile]参数使得能够插入代码
  \frametitle{归档工具}

  Java归档工具是JDK中提供的一种多用途的存档及压缩工具，可以将多个文件或
  目录合并/压缩为单个的Java归档文件（jar, java archive）。

  \pause
  
  \tta{jar文件的主要作用}

  \begin{itemize}\kai
  \item 发布和使用类库
  \item 作为程序组件或者插件程序的基本部署单位
  \item 用于打包与组件相关联的资源文件
  \end{itemize}

  \pause
  
  \tta{使用jar工具基本语法格式}

  \begin{shCode}
    >jar {-ctxui} [vfm0Me] [jar-file] [manifest-file] \  
    [entry-point] [-C dir] files ...
  \end{shCode}
\end{frame}

\begin{frame}[fragile] % [fragile]参数使得能够插入代码
  \frametitle{参数说明}

  \begin{description}\kai
  \item[-c] 创建新的归档文件。
  \item[-t] 列出归档目录。
  \item[-x] 解压缩已归档的指定（或者所有）文件。
  \item[-u] 更新现有的归档文件。
  \item[-v] 在标准输出中生成详细输出。
  \item[-f] 指定归档文件名。
  \item[-m] 包含指定清单文件中的清单信息。
  \item[-e] 为捆绑到可执行jar文件的独立应用程序指定应用程序入口点。
  \item[-0] 仅存储，不使用任何ZIP压缩。
  \item[-M] 不创建条目的清单文件。
  \item[-i] 为指定的jar文件生成索引信息。
  \item[-C] 更改为指定的目录并包含其中的文件。
  \end{description}
\end{frame}

\begin{frame}[fragile] % [fragile]参数使得能够插入代码
  \frametitle{制作并使用自己的jar文件}

  \samp{A.java}
  
  \begin{javaCode}
    public class A {
      public void ma() {
        System.out.println("In class A!");
      }
    }  
  \end{javaCode}

  \samp{TestJar.java}
  
  \begin{javaCode}
    public class TestJar {

      public static void main(String[] args) {
        A a = new A();
        a.ma();
      }
    }
  \end{javaCode}
\end{frame}

\begin{frame}[fragile] % [fragile]参数使得能够插入代码
  \frametitle{制作并使用自己的jar文件}

  \ding{182} 编译源文件A.java得到字节码文件A.class，在A.class所在路径下，
  运行如下命令进行归档处理：

  \begin{shCode}
    >jar -cvf mylib.jar *.class  
  \end{shCode}

  输出如下：

  \begin{stdoutCode}
    jar -cvf mylib.jar *.class
    added manifest
    adding: A.class(in = 380) (out= 275)(deflated 27%)  
  \end{stdoutCode}

  \ding{183} 要使用mylib.jar文件中的字节码文件，必须先将其加入到编译和
  运行环境的CLASSPATH中（注意必须指定到.jar文件的文件名）。

  \begin{shCode}
    >export CLASSPATH=".:/Users/xiaodong/temp/mylib.jar"
  \end{shCode}

  \ding{184} 编译TestJar.java源程序，并运行。
\end{frame}

\begin{frame}[fragile] % [fragile]参数使得能够插入代码
  \frametitle{发布Java应用程序}

  我们一般使用java <应用程序名字>的方式运行Java程序。学习了归档工具后，
  有了一个新的选择：

  {\Blue\hei 以归档文件的形式发布Java程序并直接从归档文件中运行。}
  
\end{frame}

\begin{frame}[fragile] % [fragile]参数使得能够插入代码
  \frametitle{发布Java应用程序}

  \samp{TestApp01.java}
  
  \begin{javaCode}
    public class TestApp01 {
      public static void main(String[] args) {
        System.out.println("App01 is running...");
      }
    }  
  \end{javaCode}

  \samp{TestApp02.java}

  \begin{javaCode}
    import java.awt.*;
    import java.awt.event.*;

    public class TestApp02 {
      public static void main(String[] args) {
        Frame f = new Frame("Test App 02");
        f.setSize(200, 200);
        f.addWindowListener(new WindowAdapter() {
          public void windowClosing(WindowEvent e) {
            System.exit(0);
          }
        });
        f.setVisible(true);
      }
    }
  \end{javaCode}
\end{frame}

\begin{frame}[fragile] % [fragile]参数使得能够插入代码
  \frametitle{发布Java应用程序过程}

  \tta{编译程序}

  \tta{程序归档发布}
  \begin{shCode}
    >jar -cfe mylib01.jar TestApp01 *.class
    >jar -cfe mylib02.jar TestApp02 *.class
  \end{shCode}

  通过使用-e参数指定当前归档文件的应用程序入口点（Entry-Point）。我们查看jar包中的清单文件可以发现多了一条Main-Class属性。
  \tta{运行程序}
  \begin{shCode}
    >java -jar mylib01.jar
    >java -jar mylib02.jar
  \end{shCode}
\end{frame}

\begin{frame}[fragile] % [fragile]参数使得能够插入代码
  \frametitle{清单文件}

  清单文件提供了归档文件的有关说明信息。jar包中使用一个特定的目录（META-INF）存放MANIFEST.MF清单文件。

  \tta{清单文件格式}

\begin{verbatim}
<属性名>:<属性值>
\end{verbatim}

  \tta{MANIFEST.MF示例}

  \begin{shCode}
    Manifest-Version: 1.0
    Created-By: 1.6.0_33 (Apple Inc.)
    Main-Class: TestApp01
  \end{shCode}

  每行最多72字符，写不下可以续行，续行必须以空格开头，且以空格开头的行都会被视为前一行的续行。
  可以自定义清单文件。
\end{frame}

\begin{frame}[fragile]
  \frametitle{本节习题}

  \wxd{简答题}
  \begin{enumerate}
  \item Java文件操作常用API有哪些？请自行搜索总结这些Java文件操作API的
    功能与那些Windows DOS命令和Linux下的Shell命令功能相
    似。Windows和Linux两个操作系统都要分别类比。举例说明：
    \begin{table}
      \footnotesize
      \setlength{\extrarowheight}{1.2mm}
      \begin{tabular}{c|c|c|c}
        \hline
        {\bf Java} & {\bf Windows DOS} & {\bf Linux Shell} & {\hei 功能说明}  \\
        \hline
        list() & dir & ls & 列出目录下文件\\
        \hline
        renameTo() & ren & mv & 修改文件名\\
        \hline
      \end{tabular}
    \end{table}
  \end{enumerate}

  \wxd{小编程}
  \begin{enumerate}
  \item 自行搜索总结Eclipse中Java Project和Maven Project导出Jar包的方法，
    并在Eclipse中编写Sample Code测试打包方法。
  \end{enumerate}

\end{frame}


% TKS %%%%%%%%%%%%%%%%%%%%%%%%%%%%%%%%%%%%%%%%%%%%%%
\begin{frame}[focus]
  \centering
  {\Huge {THE END}} \\
  \vspace{5mm}
  {\Large wangxiaodong@ouc.edu.cn} \\
\end{frame}
%%%%%%%%%%%%%%%%%%%%%%%%%%%%%%%%%%%%%%%%%%%%%%%%%%%% 
\end{document}